%% ============================================================
%% REMAINING APPENDICES -- v2 merge. None is cited by the body, so none breaks the build;
%% all are worth carrying. Source: v1 lines 2537-2650.
%%
%% NOT CARRIED OVER: v1's "BCM calibration details" (L2652-2666). It documents a global
%% x1.330 trunk multiplier and records that the dispatch model concentrated losses in the
%% terminal pipe at 4.7x nominal. Both calibrations are removed from this lineage, and the
%% second is precisely what R2.4 objected to. Reproducing the appendix would reintroduce
%% the practice the revision exists to eliminate.
%%
%% *** SEE THE NOTE AT THE END -- the per-node table below raises a substantive issue with
%% *** how the current validation table reports the temperature gates.
%% ============================================================

\section{Electricity selling price}
\label{app:selling_price}

The price received for exported electricity applies the spot price less a marketing spread,
floored, net of a self-consumption share and a fixed fee, and cannot go negative:
%
\begin{equation}
  \lambda_t^{\mathrm{sell}} = \max\!\Bigl(
    \max\bigl(\lambda_t^{\mathrm{spot}} - \lambda^{\mathrm{spread}},\;
      \lambda^{\mathrm{floor}}\bigr)
    (1 - h^{\mathrm{sell}})
    - \lambda^{\mathrm{sell,fee}},\; 0\Bigr).
\end{equation}

\section{Per-node temperature error breakdown}
\label{app:per_node_mae}

Table~\ref{tab:per_node_mae} reports supply-temperature error at all fourteen network nodes
over a 744-hour winter window, evaluated on the temperature-propagating formulation and
sorted by path length from the source. Three groups are distinguished, and the distinction
matters for how the aggregate figures in Section~\ref{subsec:validation} should be read.

\emph{Validation nodes} (six) are consumer sensors whose mixing-valve offset is below
$2$\,°C and which took no part in the calibration; their mean error is $1.32$\,°C.
\emph{Calibration nodes} (four) were used in the branch-sequential coefficient tuning, so
their residuals are in-sample by construction; their mean is $2.32$\,°C, and one node at
$3.07$\,°C would fail a $2.5$\,°C held-out criterion. That in-sample deterioration is
expected, and it is the reason a strict calibration/validation split matters.
\emph{Excluded nodes} (four) are set aside for measurement reasons unrelated to model
quality: two exhibit mixing-valve absorption above the $2$\,°C threshold and cannot be read
as junction temperatures, one is a boundary-condition reference, and one serves four
consumers with a $20.9$\,°C intra-node spread, so no single junction reference exists.

The mean across all fourteen nodes is $1.68$\,°C, between the validation-node mean of
$1.32$\,°C and the calibration-node mean of $2.32$\,°C. The validation nodes are therefore
not systematically the best-performing ones -- two of the excluded nodes are in fact the
best in the network -- so the group means are not an artefact of favourable node selection.

\medskip
{%
\captionof{table}{Per-node supply-temperature error over a 744\,h winter window, evaluated
on the temperature-propagating formulation.}
\label{tab:per_node_mae}
\centering
\footnotesize
\begin{threeparttable}
\begin{tabular*}{\columnwidth}{@{\extracolsep{\fill}} l c c c c c @{}}
\toprule
Node & Path [m] & Role & MAE [$^\circ$C] & Bias [$^\circ$C] & In gate? \\
\midrule
j$_2$    &  350 & excl. & 2.99 & $+2.98$ & No  \\
j$_3$    &  650 & excl. & 0.58 & $+0.07$ & No  \\
j$_4$    &  910 & excl. & 0.37 & $-0.18$ & No  \\
j$_9$    & 1100 & val.  & 0.87 & $+0.41$ & Yes \\
j$_5$    & 1150 & excl. & 2.35 & $+1.51$ & No  \\
j$_6$    & 1300 & cal.  & 1.62 & $-1.62$ & No  \\
j$_{10}$ & 1300 & val.  & 2.27 & $+1.78$ & Yes \\
j$_7$    & 1370 & cal.  & 3.07\tnote{a} & $+2.53$ & No  \\
j$_8$    & 1450 & cal.  & 2.29 & $+1.00$ & No  \\
j$_{11}$ & 1530 & val.  & 1.41 & $+0.71$ & Yes \\
j$_{12}$ & 1780 & val.  & 0.69 & $-0.29$ & Yes \\
j$_{13}$ & 2000 & val.  & 1.09 & $-1.02$ & Yes \\
j$_{15}$ & 2125 & val.  & 1.58 & $-0.24$ & Yes \\
j$_{14}$ & 2180 & cal.  & 2.31 & $+1.02$ & No  \\
\midrule
\multicolumn{3}{@{}l}{Group mean, validation (6 nodes)}  & 1.32 & & \\
\multicolumn{3}{@{}l}{Group mean, calibration (4 nodes)} & 2.32 & & \\
\multicolumn{3}{@{}l}{Group mean, excluded (4 nodes)}    & 1.57 & & \\
\multicolumn{3}{@{}l}{All 14 nodes}                       & 1.68 & & \\
\bottomrule
\end{tabular*}
\begin{tablenotes}\scriptsize
  \item[a] In-sample residual on a calibration node; would exceed a $2.5$\,°C held-out
    criterion.
\end{tablenotes}
\end{threeparttable}
}% end centering group

\section{Demand heterogeneity index}
\label{app:hi_definition}

With $d_n$ the share of total annual demand at node $n$, the normalised variance index is
%
\begin{equation}
  d_n = \frac{\sum_t Q_{n,t}^{\mathrm{dem}}}{\sum_n \sum_t Q_{n,t}^{\mathrm{dem}}},
  \qquad
  \mathrm{HI} = \frac{N}{N-1} \sum_{n=1}^{N}
    \left(d_n - \frac{1}{N}\right)^{2}.
  \label{eq:hi}
\end{equation}
%
The index maps monotonically to the Gini coefficient on tree topologies but compresses the
realistic range: the industrial case has $\mathrm{HI} = 0.05$ against a Gini of $0.41$. The
synthetic values $0.1$, $0.4$ and $0.8$ correspond approximately to Gini ranges of
$0.15$--$0.25$, $0.40$--$0.55$ and $0.70$--$0.85$. Because of that compression, the
model-selection observations in Table~\ref{tab:criteria} are expressed in terms of the Gini
coefficient and the top-$k$ demand share rather than this index.

%% ============================================================
%% *** SUBSTANTIVE NOTE FOR THE AUTHOR -- READ BEFORE MERGING ***
%%
%% The per-node table above is evaluated on the TEMPERATURE-PROPAGATING formulation.
%% v1 was explicit about this and about why: "The T_sup gates are assigned to [the
%% propagating variant] because only this variant propagates supply temperature along
%% pipes; [the baseline] uses fixed nominal temperatures and therefore cannot predict
%% downstream temperature deviations." (v1 L1151-1155.)
%%
%% The current tab_validation reports a far-end supply-temperature MAE of 9.21 C against a
%% 1.5 C gate, and a trunk temperature-drop MAE of 9.09 C with a SIMULATED DROP OF ZERO.
%% A simulated drop of exactly zero is the signature of a model that does not propagate
%% temperature at all. If those figures come from the node-resolved baseline, then we are
%% reporting a failed gate for a formulation that by construction cannot meet it -- which
%% is not a metering limitation and not a modelling failure, but a category error.
%%
%% This does not weaken the paper; it sharpens it, and the chain is coherent:
%%   (i)  the temperature gates apply only to a formulation that propagates temperature;
%%   (ii) the baseline does not propagate by construction, so those gates do not apply
%%        to it and should be marked "n/a" rather than "fail";
%%   (iii) the propagating level is degenerate when solved (91 % of hours at the
%%        temperature floor), so its solved output is not reportable either;
%%   (iv) temperature propagation is therefore quantified through the forward evaluator,
%%        which applies the native exponential to a fixed schedule.
%%
%% RESOLVED 2026-08-16: validation_kpis.json is the MILP baseline (boundary-condition
%% matching, nominal constant return temperature, supply temperature injected as a boundary
%% condition and not propagated). VALIDATION_METHODOLOGY.md states that the supply-
%% temperature gates are NLP-only and the MILP run skips them. tab_validation now carries a
%% Model column and marks those rows n/a for the MILP. Section 3.1 is written to match.
%%
%% NOTE: v1 had this attribution CORRECT -- its targets table assigned the supply-
%% temperature gates to the NLP and its results table scored them there. Nothing in the
%% original submission needs to be walked back on this point.
%% ============================================================
